\section{Teamwork, Leadership, and Communication (M16, M17)}
\label{sec:teamwork}

Our five-person team had distinct roles: I led software and CV, others handled hardware assembly, piloting, project administration, and flight dynamics research. In Belbin's framework~\cite{belbin2010}, I functioned as Specialist and Completer-Finisher---technical lead and de facto systems integrator. The team worked effectively where responsibilities had explicit handoff points: the hardware lead assembled the drone independently, consulting me only for Pi-to-Cube wiring; the pilot practised RC and understood abort procedures without engaging with code; the project manager booked field slots and submitted risk assessments I would have neglected.

The most productive sub-team pairing was with Edward on computer vision. His optics knowledge complemented my software skills---he calibrated camera FOV and evaluated lens distortion while I built the detection pipeline. This division by complementary expertise produced faster results than either of us working alone.

In practice, the broader software was written by one person. This was not planned---it evolved because I had the strongest programming background and found it faster to build than onboard. Reflecting through Kolb's cycle~\cite{kolb1984}, I had optimised for output over team development. I attempted to lower the barrier with numbered test scripts and documentation, but the combined context of MAVLink, state machines, and TFLite was too high for short-term engagement. Teammates expressed frustration at feeling sidelined---unable to contribute meaningfully to what became the project's largest workstream. This was the most persistent and uncomfortable tension in the team, and I bear primary responsibility for it. The lesson: contribution points must be designed in from the start, not retrofitted.

\textbf{Disagreements and compromise.} The first technical disagreement---ROS~2 versus raw MAVLink---I resolved by building a 20-line working heartbeat demo. In Tuckman's model~\cite{tuckman1965}, this storming phase was brief because the team deferred to demonstrated results. However, reflecting honestly, this was persuasion by fait accompli, not negotiation. The second disagreement taught me more. I had designed a state machine with 11 states and multiple sub-modes; the team pushed back, arguing it was too complex for anyone else to debug or modify under pressure. My instinct was to defend the design---I had tested it extensively in simulation. But after sitting with their feedback, I recognised they were right about the maintainability risk. I simplified the state transitions, removed two experimental sub-modes, and wrote explicit transition diagrams that the pilot and project manager could follow. This compromise cost me a week of refactoring, but it produced a system the whole team could reason about on flight day. It was the first time I genuinely changed a technical direction because of team input rather than technical evidence, and it shifted my view of ``good architecture'' from ``what I can debug'' to ``what the team can debug.''

\textbf{Receiving and acting on feedback.} The pilot identified that abort and confirm-target buttons in the ground station were dangerously close---a genuine safety risk I had missed. I immediately separated the controls and added a confirmation step. Beyond UI, the project manager told me my documentation was too dense---written for someone with my context, not theirs. I restructured the key guides around numbered steps and screenshots rather than architectural explanations, and created a one-page colleague checklist. Edward's feedback was more subtle: he noted that my pace of iteration made it hard for him to know which version of a module was current. I responded by adding explicit version tags to model files and maintaining a living ground-truth document that always reflected the latest state. Following Gibbs~\cite{gibbs1988}, these exchanges taught me that designing for external feedback---and acting on it visibly---is itself an engineering discipline.

\textbf{Communication across skill gaps.} Communicating the system to non-software teammates was the hardest interpersonal challenge. The \textbf{simulator} proved the most effective device: teammates understood the system intuitively when they saw the drone flying, detecting, and descending on screen. The \textbf{ground station} translated complex state into large buttons usable outdoors. \textbf{Live demonstrations} at PDR and FDR generated more engagement than slides. For the pilot specifically, I ran through the abort sequence on a bench test until he could describe what each mode indicator meant---not correcting his understanding, but developing his confidence to make independent decisions during flight. This was the closest I came to genuine team development, and I wish I had invested more in it earlier.

\textbf{Team effectiveness evaluation.} Collectively, the team delivered a working system, but the five-person structure was suboptimal for a software-heavy project. Three of the five roles (administration, flight dynamics research, piloting) had low sustained workload, while software demanded 25--30 hours per week from me alone. A more effective allocation would have paired two people on software from day one, with explicit interface contracts and code review cycles. The team's diverse backgrounds---my software depth, Edward's optics, the pilot's operational instinct---added genuine value at integration points, but the gaps between those points left people underutilised. If I were to redesign the team, I would define fewer, broader roles with more overlap, ensuring no one depends entirely on one person's availability.
